% ============================================================
%  Conclusion et Perspectives — Rapport de stage
%  Pipeline de traitement des logs LARSI / Université de Lomé
%  À insérer avant les Annexes.
% ============================================================

\documentclass[a4paper,12pt]{article}
\usepackage[utf8]{inputenc}
\usepackage[T1]{fontenc}
\usepackage[french]{babel}
\usepackage{amsmath, amssymb}
\usepackage{geometry}
\geometry{margin=2.5cm}

\begin{document}

% ------------------------------------------------------------
\chapter*{Conclusion et Perspectives}
% ------------------------------------------------------------

\section*{Conclusion générale}

Le stage effectué au sein du Laboratoire de Recherche en Sciences de l'Ingénieur (LARSI) de l'Université de Lomé avait pour mission de lever un verrou technique majeur : l'exploitation des fichiers journaux bruts, massifs et hétérogènes produits par les serveurs web Apache de l'institution. En effet, la nature textuelle et non structurée de ces journaux empêchait toute injection directe dans des pipelines d'analyse avancée ou des modèles d'apprentissage automatique. L'objectif central consistait donc à transformer ces données chaotiques en un actif analytique propre, standardisé et exploitable.

Au terme de ce travail, cet objectif a été atteint grâce à la conception et à l'implémentation d'un pipeline complet de type ETL (\emph{Extract, Transform, Load}), déployé sur l'écosystème Python (Pandas, NumPy, expressions régulières). La chaîne de traitement s'articule autour de trois modules complémentaires. Le premier, \texttt{LogParser}, assure l'ingestion récursive et le parsing orienté objet des journaux : sa capacité à détecter automatiquement et à traiter plusieurs formats (Apache Access, Apache Error, Syslog) garantit la généricité de la solution face à la diversité des sources identifiées au laboratoire (API, site web, SSL, proxy, E-learning, résultats, etc.). Le second module, \texttt{log\_cleaner}, met en œuvre un nettoyage rigoureux : imputation des valeurs manquantes par une catégorie dédiée, normalisation temporelle selon la norme ISO 8601 calée sur le fuseau UTC et déduplication stricte, qui a permis de ramener la base consolidée de 42~899 à \textbf{41~433 lignes} propres et cohérentes. Le troisième module, \texttt{log\_visualizer}, restitue graphiquement les tendances comportementales et les indicateurs structurels issus de l'analyse.

L'ingénierie des caractéristiques a constitué le cœur à forte valeur ajoutée du pipeline. À partir des champs bruts, plus de cinquante variables descriptives et comportementales ont été construites : variables temporelles (encodage cyclique sinusoïdal et cosinusoïdal de l'heure, temps normalisé), variables liées à l'IP (volume de requêtes, étalement temporel, taux de requêtes, régularité des intervalles, activité de sous-réseau, diversité des méthodes), descripteurs d'URL et de référent (profondeur, entropie de Shannon, caractères spéciaux, paramètres de requête) ainsi que des indicateurs du User-Agent (longueur, entropie, détection de bots). Ces variables ont ensuite été rendues exploitables par les algorithmes : encodage par rareté des variables catégorielles (\texttt{method\_encoded}, \texttt{status\_code\_category\_encoded}), transformation logarithmique $\log(1+x)$ des distributions asymétriques et standardisation centrée-réduite. Le jeu de données final compte ainsi \textbf{62 colonnes}, dont 57 variables numériques standardisées, prêtes pour la modélisation.

L'exploration visuelle et statistique a confirmé la pertinence de la démarche tout en révélant des caractéristiques structurantes du trafic du LARSI : une fenêtre d'observation homogène (juillet--août 2021), un trafic dominé par les requêtes GET (89,3~\%) avec une part notable de bots d'indexation (37,8~\% des requêtes), des codes de statut exclusivement en redirection (3xx) et une redondance importante entre certains descripteurs, formalisée par le clustering hiérarchique et l'analyse en composantes principales (20 composantes pour expliquer 99~\% de la variance). Ces constats guident directement les choix ultérieurs de sélection de variables et d'algorithmes.

Sur le plan professionnel, ce stage a permis de développer des compétences concrètes en ingénierie des données et en programmation Python, mais aussi d'acquérir une compréhension des enjeux de la cybersécurité, de l'audit des accès et de la détection d'anomalies. Il a également démontré que la qualité d'un projet de science des données repose en grande partie sur la rigueur de sa phase de préparation, en amont de toute modélisation.

% ------------------------------------------------------------
\section*{Perspectives}

Le pipeline réalisé constitue une fondation solide et reproductible. Plusieurs axes d'amélioration et d'extension peuvent être envisagés afin de poursuivre et de renforcer ce travail.

\subsection*{Modélisation pour la détection d'anomalies}

La suite logique du projet est l'application d'algorithmes non supervisés de détection d'anomalies sur le jeu de données standardisé : \textbf{Isolation Forest}, \textbf{PCA} (approche par erreur de reconstruction), \textbf{k-means}, \textbf{DBSCAN} ou encore les \textbf{autoencodeurs}. Ces approches permettraient d'attribuer un score d'anomalie à chaque requête, d'identifier les comportements suspects (scans, bursts de requêtes, injections) et de comparer objectivement leurs performances, en s'appuyant sur la structure de corrélation et la variance cumulée déjà analysées.

\subsection*{Enrichissement du jeu de données}

Plusieurs extensions des données sont envisageables pour affiner la détection. L'intégration d'une base \textbf{GeoIP} permettrait d'introduire un encodage par pays (\texttt{country\_encoded}) afin de signaler les origines géographiques suspectes du trafic, comme évoqué dans le chapitre consacré à l'encodage. L'exploitation des journaux d'erreurs Apache et des Syslog, déjà pris en charge par le parseur, enrichirait l'analyse d'un volet \emph{orienté serveur} complémentaire au volet \emph{orienté requête} actuel.

\subsection*{Détection en temps réel et déploiement}

Le pipeline fonctionne actuellement en traitement par lots sur des fichiers historiques. Une évolution majeure consisterait à passer à un \textbf{traitement en continu} (\emph{streaming}) : ingestion des logs au fil de l'eau, mise à jour incrémentale des statistiques de référence et alerte automatique lorsqu'un indicateur dépasse un seuil. La mise en place d'un tableau de bord de surveillance et d'un système de notification permettrait d'assurer une veille opérationnelle permanente de l'infrastructure du laboratoire.

\subsection*{Évaluation supervisée et étiquetage}

Bien que l'approche retenue soit non supervisée, la construction d'un jeu de données étiqueté (requêtes normales versus malveillantes), par validation experte ou étiquetage semi-automatique guidé par les scores d'anomalie, permettrait d'entraîner des modèles supervisés (forêts aléatoires, gradient boosting) et de mesurer l'apport de l'ingénierie des caractéristiques à l'aide de métriques objectives telles que la précision, le rappel et le score F1.

\subsection*{Amélioration continue de l'ingénierie des caractéristiques}

L'ingénierie pourrait encore être enrichie : analyse de séquences (n-grammes d'URL), représentations vectorielles (\emph{embeddings}) des User-Agents et des URLs, agrégations sur des fenêtres glissantes pour capturer des comportements localisés dans le temps, ou encore évaluation systématique de l'importance des variables pour conserver les plus discriminantes.

\subsection*{Interprétabilité et exploitation}

Enfin, l'adoption d'un tel système par les équipes de sécurité suppose une certaine \textbf{interprétabilité}. L'utilisation de techniques d'explicabilité comme SHAP ou LIME permettrait d'expliquer, pour chaque requête signalée, quelles caractéristiques ont déclenché l'alerte, facilitant ainsi la qualification des alertes et la réduction des faux positifs.

En définitive, les bases posées durant ce stage ouvrent la voie à un système complet de surveillance et de détection d'intrusions, évolutif et exploitable en conditions réelles, au croisement de l'ingénierie des données, de la cybersécurité et de l'intelligence artificielle. Ce projet constitue également un socle méthodologique précieux pour les futurs travaux de recherche du laboratoire sur l'audit des accès et la valorisation des données de l'université.

\end{document}
